%% kapa: exam=1 round=28 year=2017 period=2 form=A subject=appraisal_law subject_ko=감정평가관계법규 seq=24 number=24
%% source: Q-Net 감정평가사 기출문제 (cst00309) artlSeq=5199603 file=extracted/1008/5199603/2017년 제28회 감정평가사 1차시험 문제지/2교시 A형.pdf
%% conversion: kordoc 4.24.51 -> tools/gen1.py (auto); review=proofread-sonnet; flags=none
\begin{question}{24}
건축법령상 건축허가 등에 관한 설명으로 옳지 않은 것은?
\begin{choices}
  \item 광역시에 연면적의 합계가 20만 제곱미터인 공장을 건축하려면 광역시장의 허가를 받아야 한다.
  \item 허가권자는 숙박시설에 해당하는 건축물의 용도가 교육환경 등 주변 환경을 고려할 때 부적합하다고 인정되는 경우에는 건축위원회의 심의를 거쳐 해당 건축허가를 하지 아니할 수 있다.
  \item 건축허가를 받으면「도로법」에 따른 도로의 점용 허가를 받은 것으로 본다.
  \item 건축 관련 입지와 규모에 대한 사전결정을 신청한 자는 사전결정을 통지받은 날부터 2년 이내에 건축허가를 신청하여야 한다.
  \item 건축허가를 받은 후 건축주를 변경하는 경우에는 신고하여야 한다.
\end{choices}
\end{question}
\answer{1}
